\documentclass[11pt]{article}

\usepackage[a4paper,margin=1in]{geometry}
\usepackage{amsmath,amssymb,mathtools}
\usepackage{booktabs,tabularx,array}
\usepackage[hidelinks]{hyperref}

\newcommand{\R}{\mathbb{R}}
\newcommand{\eps}{\varepsilon}
\newcommand{\Sym}{\operatorname{Sym}}

\title{Notation and Terminology for Polynomial-Zonotope Neural-Network Certification}
\author{}
\date{}

\begin{document}
\maketitle

This note fixes the notation used in the accompanying references.  There are
exactly two classes of noise symbols: domain noise symbols \(\alpha\) and
approximation noise symbols \(\eta\).  Their combined vector is
\(\eps=(\alpha,\eta)\).  The word \emph{residual} is reserved for the actual
approximation difference, not used as a second name for \(\eta\).

\section{Basic objects}

\renewcommand{\arraystretch}{1.25}
\small
\noindent\begin{tabularx}{\textwidth}{@{}>{$}l<{$} X@{}}
\toprule
\text{Symbol} & \text{Meaning} \\
\midrule
x & Physical input point in the domain of the neural network. \\
\mathcal X\subseteq\R^{d_0} & Physical input set. \\
X(\alpha) & Polynomial-zonotope parametrization of \(\mathcal X\). \\
\alpha=(\alpha_1,\ldots,\alpha_p) & Domain noise symbols, with \(\alpha\in[-1,1]^p\). \\
\eta=(\eta_1,\ldots,\eta_q) & Approximation noise symbols introduced by certified activation enclosures, with \(\eta\in[-1,1]^q\). \\
\eps=(\alpha,\eta) & Combined vector of all polynomial variables after propagation. \\
\lambda & Generic multi-index in \(\mathbb N_0^{p+q}\). \\
\lambda=(\lambda^\alpha,\lambda^\eta) & Split into domain and approximation components. \\
\eps^\lambda=\alpha^{\lambda^\alpha}\eta^{\lambda^\eta} & Monomial associated with \(\lambda\). \\
\Lambda_P & Finite exponent support of a sparse polynomial \(P\). \\
\bottomrule
\end{tabularx}

The input polynomial zonotope and a general \(U\)-valued propagated
polynomial are written as
\[
 X(\alpha)=x_0+\sum_{\lambda\in\Lambda_X}x_\lambda\alpha^\lambda,
 \qquad
 P(\eps)=\sum_{\lambda\in\Lambda_P}P_\lambda\eps^\lambda,
 \quad P_\lambda\in U.
\]
When several multi-indices occur in one formula, auxiliary letters such as
\(\kappa\) and \(\gamma\) may be used; \(\lambda\) remains the canonical
generic multi-index.

\section{Network and two-jet notation}

\noindent\begin{tabularx}{\textwidth}{@{}>{$}l<{$} X@{}}
\toprule
\text{Symbol} & \text{Meaning} \\
\midrule
\Phi & Complete neural network. \\
A_\ell(y)=W_\ell y+b_\ell & Affine layer, with weight matrix \(W_\ell\) and bias \(b_\ell\). \\
\sigma & Scalar activation function. \\
\sigma_\ell & Componentwise activation layer. \\
z_\ell,\ y_\ell & Preactivation and postactivation at hidden layer \(\ell\). \\
Z_\ell(\eps),\ Y_\ell(\eps) & Corresponding polynomial enclosures. \\
J_\ell(\eps),\ H_\ell(\eps) & Polynomial enclosures of the first and second physical-input derivatives. \\
\left.\partial F\right|_x & First derivative of \(F\) evaluated at \(x\). \\
\left.\partial^2F\right|_x & Second derivative of \(F\) evaluated at \(x\). \\
\left.\partial_aF^i\right|_x,\ \left.\partial_{ab}F^i\right|_x & Coordinate entries of these derivatives. \\
J_X(\alpha)=\left|\det\left.\partial X\right|_\alpha\right| & Nonnegative geometric Jacobian density of the domain parametrization. \\
\bottomrule
\end{tabularx}

We use
\[
 \Phi=A_L\circ\sigma_L\circ\cdots\circ A_1\circ\sigma_1\circ A_0,
\]
and the propagated output coefficients have types
\[
 Y_\lambda\in\R^{d_{\rm out}},\qquad
 J_\lambda\in\R^{d_{\rm out}}\otimes\R^{d_0*},\qquad
 H_\lambda\in\R^{d_{\rm out}}\otimes\Sym^2(\R^{d_0*}).
\]
The derivatives are always with respect to the physical input \(x\); the
dependence on \(\alpha\) only records evaluation along \(X(\alpha)\).

\section{Activation enclosures and terminology}

For \(r\in\{0,1,2\}\), with \(\sigma^{(0)}=\sigma\), the standard enclosure is
\[
 \sigma^{(r)}(t)
 \in \widehat\sigma_{\ell i}^{(r)}(t)
     +\rho_{\ell i}^{(r)}\eta_{\ell i}^{(r)},
 \qquad \eta_{\ell i}^{(r)}\in[-1,1].
\]

\noindent\begin{tabularx}{\textwidth}{@{}>{$}l<{$} X@{}}
\toprule
\text{Symbol or term} & \text{Meaning} \\
\midrule
\widehat\sigma^{(r)} & Polynomial or affine approximation of \(\sigma^{(r)}\) on the certified preactivation interval. \\
\rho\geq0 & Certified approximation-error radius. \\
\eta^{(r)} & Approximation noise symbol. \\
\rho\eta & Approximation noise term; indices are added when needed. \\
r(t) & Residual function, defined by
\(r(t)=\sigma^{(r)}(t)-\widehat\sigma^{(r)}(t)\). \\
\bottomrule
\end{tabularx}
\normalsize

A fresh approximation noise symbol is introduced for every independently
certified scalar enclosure.  In the general two-jet construction this gives
\(\eta_{\ell i}^{(0)},\eta_{\ell i}^{(1)},\eta_{\ell i}^{(2)}\) per neuron.
The same \(\eta_{\ell i}^{(1)}\) is reused wherever the same enclosure of
\(\sigma'\) occurs, including both the Jacobian and Hessian recurrences.
Special constructions may instead derive all activation derivatives from one
shared enclosure; such dependency sharing must be stated explicitly.

\section{Canonical vocabulary}

\noindent\begin{tabularx}{\textwidth}{@{}l X@{}}
\toprule
Preferred term & Usage \\
\midrule
Domain noise symbol & A component of \(\alpha\). \\
Approximation noise symbol & A component of \(\eta\); this is the canonical name. \\
Approximation-error radius & The nonnegative coefficient \(\rho\). \\
Approximation noise term & The product \(\rho\eta\). \\
Residual function & The actual difference between a function and its approximation. \\
Polynomial enclosure & A polynomial together with its approximation noise terms. \\
\bottomrule
\end{tabularx}

The terms \emph{residual symbol}, \emph{remainder symbol}, and
\emph{pointwise-error variable} are avoided as alternative names for \(\eta\).
The approximation noise symbol is a formal enclosure variable; its coefficient
is the certified uniform bound on the residual over the relevant interval.

\end{document}
